\documentclass[11pt]{article}
\usepackage[a4paper,margin=1in]{geometry}
\usepackage{amsmath,amssymb,amsthm}
\usepackage{booktabs}
\usepackage{graphicx}
\usepackage{xcolor}
\usepackage{hyperref}
\hypersetup{colorlinks=true,linkcolor=blue,urlcolor=blue,citecolor=blue}
\newcommand{\rt}[1]{\texttt{#1}}
\newcommand{\fn}[1]{\texttt{#1}}
\newcommand{\att}{\mathrm{att}}
\newcommand{\wcc}{\mathrm{wcc}}
\newcommand{\ifgraphic}[2]{\IfFileExists{#1}{\includegraphics[width=#2]{#1}}%
  {\fbox{\parbox{#2}{\centering\small figure \texttt{\detokenize{#1}} pending}}}}
\newcommand{\iftable}[1]{\IfFileExists{#1}{\input{#1}}%
  {\fbox{\parbox{0.9\textwidth}{\centering\small table
  \texttt{\detokenize{#1}} pending}}}}

\title{\textbf{Report R15 --- the tower in the attractor plane, and whether the
resets are idle inside it}\\
\large Krylov sector analysis of quantum cellular automata}
\author{Tier 1 analysis}
\date{2026-08-05}

\begin{document}
\maketitle

\begin{abstract}
In the right-hand panel of R9's F1 a column of V+reset rules stands at
$a_\att=1$ with $b_\att$ running from $1.41$ to $1.89$: a bounded number of
attractors, the largest of which still grows exponentially. There are $68$ such
rules of the $160$ V+reset ones. Their growth bases are not spread out --- $61$
of the $68$ sit on five \emph{exact} algebraic constants, $\varphi$ and $\sqrt3$
taking $21$ each. The natural explanation is that the dissipative channels never
act inside the largest attractor, so what grows there is really the underlying
unitary. \textbf{That is true for exactly $12$ of the $68$, and false for the
other $56$.} For the twelve it is true in the strongest sense: the largest
attractor \emph{is} the largest sector of the coherent parent, state for state,
and all twelve sit at $\varphi$. For the other fifty-six the resets do work at
$\Theta(N)$ sites --- the count of firing sites grows linearly with system size,
so this is not a boundary effect --- and the exponential growth survives
dissipation rather than avoiding it. Nine rules reach $\varphi$ \emph{with}
firing resets, so the constant does not determine the mechanism. \rt{obc0},
$N=10$ for the census.
\end{abstract}

\section{The tower}

A rule is in the tower if it is V+reset, its attractor \emph{count} is flat
($a_\att=1$ to within $0.02$) and its largest attractor grows exponentially
($b_\att>1.02$). That is $68$ rules.

\begin{figure}[htbp]
\centering
\ifgraphic{../../figures/fig_tower_obc0.pdf}{\textwidth}
\caption{\textbf{Left:} the tower, one bar per distinct $b_\att$, gold marking
the rules on which no reset ever fires. \textbf{Right:} the hypothesis tested
--- the number of sites at which some reset does work on the largest attractor,
against $N$. Flat at zero for twelve rules; linear in $N$ for the rest, tracking
the dotted ``all $N$ sites'' line.}
\end{figure}

\subsection{Where the bases sit}

\begin{center}
\iftable{tab_r15_bases_obc0.tex}
\end{center}

\noindent Five constants absorb $61$ of the $68$, and for those the base comes
from an exact integer recurrence rather than a fit:
$\sqrt2$ ($3$ rules), $\varphi$ ($21$), $\sqrt3$ ($20$ of $21$),
$2\cos(\pi/8)=\sqrt{2+\sqrt2}$ ($8$) and $\sqrt{(3+\sqrt{17})/2}$ ($9$). The
remaining seven are fits. Two of them, at $1.847632$, are within
$1.3\times10^{-4}$ of $2\cos(\pi/8)$ and are almost certainly that constant
imperfectly fitted; the pair at $1.5537/1.5539$ and the pair at $1.7180$ have no
identification here and are left as fits rather than rounded into a name.

These are the constants of \emph{constrained binary words} --- Perron roots of
small transfer matrices --- which is exactly the company the unitary rules keep.
That is what makes the idle-reset hypothesis so natural.

\section{The hypothesis, tested}

\begin{quote}
\emph{The dissipative channels never fire on the largest attractor, so what
grows there is the underlying unitary.}
\end{quote}

\noindent A reset \rt{D} writes $0$ and \rt{E} writes $1$; either does
\emph{work} only when the target bit is not already at that value, and
otherwise acts as the identity on that state. \fn{tower.firing\_sites} runs the
brick-wall sweep by hand, following every Hadamard branch, and returns the sites
at which some branch had a reset do work. (The hand sweep is validated against
the engine's successor map in the tests: it must reproduce \fn{wcc.make\_succ}
exactly, over all states of $37$ rules at three sizes, or the firing flags mean
nothing.)

\begin{center}\small
\begin{tabular}{lrl}
\toprule
& rules & \\
\midrule
no reset ever fires on the largest attractor & $\mathbf{12}$ & hypothesis holds \\
some reset fires & $\mathbf{56}$ & hypothesis fails \\
\bottomrule
\end{tabular}
\end{center}

\subsection{The twelve, and why they work}

\begin{center}
\iftable{tab_r15_idle_obc0.tex}
\end{center}

\noindent For all twelve the statement holds in the strongest available form:
the child's successor map \emph{coincides} with the coherent parent's on every
state of the largest attractor, and the attractor has exactly the size of the
parent's largest sector. It is not merely that the resets happen to be idle ---
the attractor \emph{is} a sector of the parent, state for state.

They fall into two families of six, and the mechanism is visible in the tuple:

\begin{center}\small
\begin{tabular}{llll}
\toprule
family & tuples & parent & $|A|$ at $N=10$ \\
\midrule
\rt{V} at $r_{00}$, rest from $\{$\rt{D},\rt{I}$\}$ &
  \rt{VDDD}, \rt{VIDD}, \rt{VDID}, \rt{VDDI}, \rt{VIDI}, \rt{VDII} &
  $W201$ & $144=F_{12}$ \\
\rt{V} at $r_{11}$, rest from $\{$\rt{E},\rt{I}$\}$ &
  \rt{IEIV}, \rt{EEIV}, \rt{IIEV}, \rt{EIEV}, \rt{IEEV}, \rt{EEEV} &
  $W108$ & $55=F_{10}$ \\
\bottomrule
\end{tabular}
\end{center}

\noindent The reset's target value agrees with what those sites already carry.
In the first family the Hadamard fires only on the neighbourhood $00$, and every
\rt{D} sits at a neighbourhood that occurs only when the site is already $0$, so
writing $0$ changes nothing. The second family is the same statement with
$0\leftrightarrow1$ and \rt{E} for \rt{D}. The surviving set is the parent's
Fibonacci-counted sector, hence $\varphi$ --- and indeed \textbf{every one of the
twelve sits at $\varphi$}, none at any other constant.

\subsection{The fifty-six, and why the hypothesis fails for them}

For these the resets fire, and not marginally: the number of \emph{sites} at
which some reset does work on the largest attractor grows linearly with $N$
(right panel of the figure), typically reaching $N-1$ or $N$. So it cannot be
rescued as a boundary effect or as a bounded set of defect sites --- dissipation
acts throughout the bulk of the attractor, and the exponential growth survives
it rather than avoiding it.

What these rules share with the twelve is only the \emph{form} of the answer:
an algebraic Perron root. A reset that fires deterministically given its
neighbourhood is still a local constraint, and a set of words cut out by local
constraints is still counted by a transfer matrix. The constant does not
require the dynamics to be unitary on the attractor --- only local.

\subsection{The constant does not identify the mechanism}

Twenty-one rules sit at $\varphi$; twelve are idle and \textbf{nine are not}
--- rules $11$, $47$, $81$, $117$, $139$, $161$, $171$, $209$ and $241$ reach
the same golden ratio with resets firing at $\Theta(N)$ sites. So one cannot
read the mechanism off the growth base, which is worth saying because the
temptation to do so is exactly what this report set out to test.

\section{What this settles and what it does not}

\begin{itemize}
\item \textbf{Settled.} The idle-reset explanation is a real mechanism and it is
  a minority one: $12$ of $68$, all at $\varphi$, all inheriting a parent sector
  outright.
\item \textbf{Settled.} For the other $56$ the resets act on $\Theta(N)$ sites of
  the attractor, so no version of ``the channel is secretly unitary there''
  survives.
\item \textbf{Open.} The transfer matrices themselves. $\sqrt3$, $2\cos(\pi/8)$
  and $\sqrt{(3+\sqrt{17})/2}$ each take $8$--$21$ rules, and the constraint
  each one counts has not been identified here; that would need the wall grammar
  of Tier 2 applied to the recurrent set rather than to the sector.
\item \textbf{Open.} The two unidentified pairs, at $1.5537$ and $1.7180$.
\item \textbf{Caveat.} \rt{obc0}, and the census is at $N=10$ with the
  firing-site scaling checked at $N=8,10,12$. Attractors are monitored-only
  objects, so everything here is a statement about the measured chain.
\end{itemize}

\section*{Reproduce}

\begin{verbatim}
python -m qca_fragmentation.viz.tower --bc obc0 --N 10
\end{verbatim}

\noindent writes \fn{fig\_tower\_obc0} and the two tables. The module is
\fn{code/qca\_fragmentation/viz/tower.py}; \fn{idle\_census(rule, N)} returns the
per-rule verdict and \fn{firing\_sites(x, N, t)} the raw site set for one state.

\end{document}
